% Compiler: pdfLaTeX
\documentclass[11pt,a4paper]{article}

\usepackage[utf8]{inputenc}
\usepackage[T1]{fontenc}
\usepackage{lmodern}
\usepackage{microtype}
\usepackage{geometry}
\usepackage{enumitem}
\usepackage{titlesec}
\usepackage{fancyhdr}
\usepackage{hyperref}

% Document details
\newcommand{\researchtitle}{The tragedy of the commons}
\newcommand{\researchshorttitle}{The commons}
\newcommand{\researchdate}{2026-10-25}

% Layout
\geometry{left=2.7cm,right=2.7cm,top=2.5cm,bottom=2.5cm,headheight=14pt}
\hypersetup{hidelinks,pdftitle={\researchtitle}}
\urlstyle{same}
\frenchspacing
\setlength{\parindent}{0pt}
\setlength{\parskip}{0.55em}
\linespread{1.04}
\setlength{\emergencystretch}{1em}

\titleformat{\section}
    {\large\bfseries}{\thesection}{0.65em}{}
\titleformat{\subsection}
    {\normalsize\bfseries}{\thesubsection}{0.65em}{}
\titlespacing*{\section}{0pt}{1.15em}{0.35em}
\titlespacing*{\subsection}{0pt}{0.7em}{0.2em}
\setlist[itemize]{leftmargin=1.35em,itemsep=0.25em,topsep=0.2em}
\setlist[enumerate]{leftmargin=1.6em,itemsep=0.25em,topsep=0.2em}

\pagestyle{fancy}
\fancyhf{}
\fancyhead[L]{\small Random Research}
\fancyhead[R]{\small\researchshorttitle}
\fancyfoot[C]{\small\thepage}
\renewcommand{\headrulewidth}{0pt}
\renewcommand{\footrulewidth}{0pt}

\begin{document}

{\LARGE\bfseries\researchtitle\par}
\vspace{0.4em}
{\small\researchdate\par}

\section{Overview}

A shared resource can be overused when each person receives the benefit
of taking more while spreading part of the damage across everyone else.
Hardin illustrated this problem with a pasture open to all: each herder
has an incentive to add animals, although the land has limits.
\cite{hardin}

The model is easy to follow. Whether it explains a particular community
depends on the rules and circumstances of that community.

\section{Findings}

Hardin's pasture is a thought experiment about incentives. The extra
animal benefits its owner, while the effects of overgrazing are shared.
Repeating that decision across users can damage the resource on which
they all depend. \cite{hardin}

Ostrom challenges the assumption that resource users cannot cooperate
without central direction. Her review draws on experiments and field
research to examine how people use norms and rules to manage shared
resources. It also treats policy changes as attempts that can fail,
requiring attention to their effects. \cite{ostrom}

In her NYU lecture, Ostrom discussed irrigation systems in Nepal.
The university's account reports that locally managed systems could
perform as well as or better than systems installed by development
agencies. This challenges a simple ranking of management arrangements;
it does not establish that local management always succeeds. \cite{nyu}

\section{Example}

Imagine several households drawing water from one well. Each household
benefits from taking enough for its own needs. If everyone increases
withdrawals during a dry period, less water remains for the group.

Now suppose the households agree on limits and a way to check use.
The shared resource is the same, but the decisions take place under
different rules. This hypothetical example explains why access and
management need to be considered separately.

\section{Discussion}

The word tragedy makes the outcome sound settled before anybody acts.
That is what I find limiting about the phrase. Hardin's example helps
explain the pressure to overuse a resource, but it leaves me wanting to
know who can make rules and whether people will follow them.

Ostrom's account feels more useful for asking those questions. It also
asks more of the reader. Saying that rules should suit local conditions
does not tell us which rules will work in a specific place.

I would be suspicious of any argument that uses the commons to justify
one form of ownership without looking at how it is actually managed.

\section{Practical implications}

For a shared resource, I would start by identifying who uses it, what
limits it has, and who bears the cost of overuse. Then I would ask how
users can agree on rules and respond when those rules fail.

The people affected need some way to see whether the arrangement is
working. Good intentions alone leave the original incentive problem
unresolved.

\renewcommand{\refname}{References}
\begin{thebibliography}{9}
\small
\setlength{\itemsep}{0.35em}

\bibitem{hardin}
Hardin, G. (1968).
\href{https://doi.org/10.1126/science.162.3859.1243}
{The tragedy of the commons.}
\textit{Science}, 162(3859), 1243--1248.

\bibitem{ostrom}
Ostrom, E. (1999).
\href{https://doi.org/10.1146/annurev.polisci.2.1.493}
{Coping with tragedies of the commons.}
\textit{Annual Review of Political Science}, 2, 493--535.

\bibitem{nyu}
NYU School of Law.
\href{https://www.law.nyu.edu/news/OSTROM_CONFERENCE_KEYNOTE}
{Nobel Prize winner Elinor Ostrom delivers keynote at conference on cultural commons.}
Report of Ostrom's lecture.

\end{thebibliography}
\end{document}
